\documentclass{article}
\usepackage[T1]{fontenc}
\usepackage{lmodern}
\usepackage{graphicx}
\usepackage{amsmath,amssymb}
\usepackage[a4paper,margin=2cm]{geometry}
\usepackage[absolute,overlay]{textpos}
\setlength{\TPHorizModule}{1mm}
\setlength{\TPVertModule}{1mm}
\usepackage[indonesian]{babel}
\usepackage{hyperref}
\setlength{\emergencystretch}{2em}
% unit: Halaman 55

\begin{document}

% === Halaman 55 (llm) ===

\textbf{Kriptografi di Eropa dari Zaman Renaisans sampai Abad 19}

\begin{itemize}
    \item Zaman renaisans $\rightarrow$ abad pertengahan (abad 15-16)
    \item \textit{Cipher} terkenal pada abad pertengahan hingga abad 19:
    \begin{enumerate}
        \item \textit{Vigenere Cipher}\
        Dipublikasikan oleh diplomat Perancis bernama Blaise de Vigenere pada tahun 1586.
        \vspace{5mm}
        \item \textit{Playfair Cipher}\
        Dipromosikan oleh diplomat Inggris, Lord Playfair, meskipun penemu aslinya adalah Charles Wheatstone pada tahun 1854.
    \end{enumerate}
\end{itemize}

\end{document}
